\scanpage{60}
2°) $\pi^\Gamma=(1)$ (i.e. il n'y a pas d'autres pt fixes dans $\widetilde U$ sur le $\unsure{=|\overline m}$ $x\in U(K)$), et enfin

3°) il n'y a pas en $\unsure{=|\overline m}$ temps de pt fixe $\unsure{\widetilde x}$ à l'infini --- i.e. il n'existe pas de $L_i$ normalisé par $\Gamma$, i.e. un scindage du premier type n'est pas en $\unsure{=|\overline m}$ temps du deuxième (fait que nous avons établi directement, précédemment).

\margin{C'est une particularisation d'une ``conséquence conjecturale'' de la conjecture du $\S1$ --- sauf qu'il faudrait écrire $\pi^{\Gamma^\circ}=(1)$}

D'autre part, dans le cas de pts fixes à l'$\infty$, l'unicité de l'image dans $X$ signifie qu'une $\unsure{=|\overline m}$ action effective ne peut normaliser à la fois un $L_i$ et un $L_j$ ($i\ne j$) --- fait également établi sauf dans le cas $(g,\nu)=(0,3)$ --- et l'unicité --- au dessus d'un $i\in\bar I$ fixé signifie que le $L_i$ \inserted{($i$ fixé)} normalisé par $\Gamma$ est unique, ce qui est un affaiblissement de la relation $L_i=\deleted{\unsure{Z_\pi|Z_{\mathrm{en}}}}\pi^{\Gamma^\circ}$ pour ces opérations, conjecturée plus haut.
\margin{59}
